\chapter{La rete}
\label{chap:the-network}

\note{in questo capitolo non mi sono venuti in mente nomi migliori per le sezioni,
con magari meno inglesismi e/o meno legati al codice, però i concetti sono questi.}

\section{YAAIR}
\label{sec:yaair}
\sidenote{essendo una sigla l'ho messa in maiuscolo, poi nel testo sarà un acronimo}

Descrizione di yaair, come funziona, il modello a strati

\subsection{Il trait Network}
\label{sec:the-network-trait}

Spiegazione più dettagliata del concetto di network e dei due metodi da cui è composta.

\section{La ZenohNetwork}
\label{sec:the-zenoh-network}
\sidenote{La generalizzazione per zenoh e zenoh pico devo ancora implementarla, quindi
potrebbe cambiare questa sezione}

Spiegazione del funzionamento della rete zenoh, con anche qualche schema.
\note{riguardo agli schemi, ho pensato di metterli in quanto mi avete detto
che dovrei seguire una struttura simile alla relazione del progetto di OOP, ma
essendo il lavoro fino a questo punto solamente un wrapper di una libreria esistente,
non c'è davvero una parte di analisi. Al massimo, giusto riguardo a come ho gestito
le closure in zenoh pico (in cui avevo anche messo la possibilità di usarle in maniera
asincrona invece che come callback, ma che poi non è servito)}

\subsection{Il MessageStore}
\label{sec:the-message-store}

Spiegazione del concetto di message store, della gestione della concorrenza e della
lifespan dei nodi in base all'ultimo update.

\subsubsection{La task di hearthbit}
\label{sec:the-heartbeat-task}

Spiegazione della hearthbit task per mantenere vivo un nodo.

\subsection{Implementazione tramite Zenoh Pico}
\label{sec:zenoh-pico-implementation}

Approfondimento sull'implementazione zenoh pico

\subsection{Implementazione tramite Zenoh}
\label{sec:zenoh-implementation}

Approfondimento sull'implementazione zenoh
